% ============================================================
%  Conclusion Générale et Perspectives
%  Titre : gras centré 24pt Times New Roman
% ============================================================
\clearpage
\thispagestyle{plain}

\vspace*{2cm}

\begin{center}
{\fontsize{24}{29}\selectfont\bfseries Conclusion Générale et Perspectives}
\end{center}
\addcontentsline{toc}{chapter}{Conclusion Générale et Perspectives}

\vspace{1.5cm}

Le projet \textbf{Dental Reservation} a permis de concevoir et réaliser une solution web complète, moderne et fonctionnelle pour la gestion des cabinets dentaires. L'application couvre l'ensemble du cycle de gestion~: de l'authentification sécurisée du praticien, à la gestion des patients et de leurs dossiers médicaux, en passant par la planification intelligente des rendez-vous avec détection automatique des conflits de créneaux.

Le choix d'une architecture client-serveur avec un backend ASP.NET Core 8 et un frontend React TypeScript a permis une séparation claire des responsabilités, facilitant la maintenance et l'évolution du système. L'utilisation d'Entity Framework Core comme ORM a simplifié les interactions avec la base de données SQL Server tout en protégeant contre les injections SQL. TanStack React Query a optimisé la gestion de l'état serveur côté frontend, offrant un cache intelligent et une synchronisation automatique.

La modélisation UML réalisée avec PlantUML (8 diagrammes) a fourni une base solide de conception avant l'implémentation, couvrant les cas d'utilisation, les classes, les séquences d'interaction et les flux d'activité.

La conteneurisation avec Docker et Docker Compose assure un déploiement reproductible et portable, et la documentation API via Swagger facilite l'intégration et les tests.

\section*{Perspectives d'évolution}

Plusieurs axes d'amélioration et d'extension sont envisagés pour les versions futures~:

\begin{enumerate}
    \item \textbf{Application mobile}~: Développement d'une application mobile React Native ou Flutter permettant aux patients de prendre rendez-vous directement depuis leur smartphone et de recevoir des rappels.

    \item \textbf{Notifications push et rappels}~: Intégration d'un service d'envoi de rappels par email (via SendGrid) ou SMS (via Twilio) avant chaque rendez-vous, réduisant le taux de rendez-vous manqués.

    \item \textbf{Gestion des paiements et facturation}~: Ajout d'un module de facturation permettant de générer des factures, suivre les paiements et s'intégrer avec des solutions de paiement en ligne.

    \item \textbf{Tableau de bord analytique avancé}~: Intégration de graphiques interactifs avec des bibliothèques comme Recharts ou Chart.js pour visualiser l'évolution des rendez-vous, le taux de remplissage, et les tendances.

    \item \textbf{Multi-cabinet et multi-rôles}~: Support de plusieurs cabinets avec des rôles différenciés (administrateur, praticien, secrétaire, patient) et des permissions granulaires.

    \item \textbf{Intégration avec des systèmes tiers}~: Connexion avec des logiciels de comptabilité, des systèmes d'information hospitaliers (SIH), ou des annuaires de professionnels de santé.

    \item \textbf{Hachage des mots de passe}~: Remplacement du stockage en clair des mots de passe par un algorithme de hachage sécurisé (bcrypt ou Argon2) pour une sécurité renforcée.

    \item \textbf{Tests automatisés}~: Mise en place d'une suite de tests unitaires (xUnit pour le backend, Vitest pour le frontend) et de tests d'intégration pour garantir la fiabilité du code.

    \item \textbf{Internationalisation (i18n)}~: Support multilingue (français, anglais, arabe) pour les cabinets opérant dans des environnements multilingues.

    \item \textbf{Prise de rendez-vous en ligne par les patients}~: Portail web dédié permettant aux patients de consulter les créneaux disponibles et de réserver eux-mêmes leurs rendez-vous.
\end{enumerate}

\clearpage
